%% kapa: exam=1 round=21 year=2010 period=2 form=A subject=accounting subject_ko=회계학 seq=36 number=36
%% source: Q-Net 감정평가사 기출문제 (cst00309) artlSeq=5024181 file=extracted/1008/5024181/2010년 제21회 감정평가사 1차시험 기출문제/2교시.hwp
%% conversion: kordoc 4.24.51 -> tools/gen1.py (auto); review=proofread-sonnet; flags=none
\begin{question}{36}
(주)한국은 종합원가계산제도를 도입하고 있다. 20×0년 1분기 동안 생산관련 자료는 다음과 같다.

{\small\centering
\begin{tabular}{lr}
\hline
당기투입량 & 8,000톤 \\
완성품 & 7,500톤 \\
기말재공품 & 2,000톤 \\
기초재공품 & 1,500톤 \\
\hline
\end{tabular}\par}

가공원가는 공정전반에 걸쳐 균등하게 발생한다. 기말재공품은 세 개의 완성도로 구성되어 있는데, 기말재공품의 1/4은 완성도가 80\%이며, 1/2은 50\%, 나머지 1/4은 20\%이다. 선입선출법(FIFO)을 적용할 경우 가공원가의 완성품 환산량이 7,960톤이라면 20×0년 1분기 기초재공품의 완성도는 얼마인가? (단, 1분기 기초재공품은 한 개의 완성도로만 구성됨)
\begin{choices}
  \item 32\%
  \item 36\%
  \item 40\%
  \item 44\%
  \item 48\%
\end{choices}
\end{question}
\answer{2}
